% ═══════════════════════════════════════════════════════════════════════════════
%  ISLS Validation Harness v1.0.0
%  Deterministic Benchmark, Test & Operational Metering Framework
%  Sebastian Klemm — 12 March 2026
% ═══════════════════════════════════════════════════════════════════════════════
\documentclass[10pt,a4paper]{article}
\usepackage[utf8]{inputenc}\usepackage[T1]{fontenc}\usepackage{lmodern}
\usepackage[english]{babel}
\usepackage{amsmath,amssymb,mathtools}
\usepackage[left=20mm,right=20mm,top=25mm,bottom=25mm]{geometry}
\usepackage{hyperref}\usepackage{enumitem}\usepackage{booktabs}
\usepackage{array,tabularx,longtable}\usepackage{xcolor}
\usepackage{listings}\usepackage{fancyhdr}\usepackage{titlesec}
\usepackage{tcolorbox}\usepackage{microtype}\usepackage{float}

\definecolor{accent}{HTML}{0369A1}\definecolor{accentdark}{HTML}{0C4A6E}
\definecolor{success}{HTML}{059669}\definecolor{warn}{HTML}{D97706}
\definecolor{danger}{HTML}{DC2626}\definecolor{codebg}{HTML}{F8FAFC}
\definecolor{codeborder}{HTML}{E2E8F0}\definecolor{defbg}{HTML}{F0FDF4}
\definecolor{defborder}{HTML}{86EFAC}\definecolor{thmbg}{HTML}{E0F2FE}
\definecolor{thmborder}{HTML}{38BDF8}\definecolor{sectioncolor}{HTML}{0C4A6E}
\definecolor{rustbg}{HTML}{FFF7ED}\definecolor{rustborder}{HTML}{EA580C}
\definecolor{harnessbg}{HTML}{FDF4FF}\definecolor{harnessborder}{HTML}{A855F7}

\hypersetup{colorlinks=true,linkcolor=accent,citecolor=accent,urlcolor=accentdark,
  pdftitle={ISLS Validation Harness v1.0.0},pdfauthor={Sebastian Klemm}}
\pagestyle{fancy}\fancyhf{}
\fancyhead[L]{\scriptsize\textcolor{gray}{ISLS Validation Harness v1.0.0}}
\fancyhead[R]{\scriptsize\textcolor{gray}{Sebastian Klemm — 12 March 2026}}
\fancyfoot[C]{\scriptsize\textcolor{gray}{\thepage}}
\renewcommand{\headrulewidth}{0.3pt}\renewcommand{\footrulewidth}{0.1pt}

\titleformat{\section}{\large\bfseries\color{sectioncolor}}{\thesection}{0.8em}{}[\vspace{-0.4em}\textcolor{accent}{\rule{\textwidth}{0.6pt}}]
\titleformat{\subsection}{\normalsize\bfseries\color{accentdark}}{\thesubsection}{0.8em}{}
\titleformat{\subsubsection}{\small\bfseries\color{accent}}{\thesubsubsection}{0.8em}{}
\titlespacing*{\section}{0pt}{1.5ex plus .5ex minus .2ex}{1ex plus .2ex}
\titlespacing*{\subsection}{0pt}{1ex plus .3ex minus .1ex}{0.5ex plus .1ex}

\tcbuselibrary{skins,breakable}
\newtcolorbox{specbox}[1][]{colback=thmbg,colframe=thmborder,arc=1.5pt,left=4pt,right=4pt,top=2pt,bottom=2pt,title={\scriptsize\bfseries #1},fonttitle=\color{accent},breakable}
\newtcolorbox{rustbox}[1][]{colback=rustbg,colframe=rustborder,arc=1.5pt,left=4pt,right=4pt,top=2pt,bottom=2pt,title={\scriptsize\bfseries #1},fonttitle=\color{rustborder!80!black},breakable}
\newtcolorbox{cmdbox}[1][]{colback=codebg,colframe=codeborder,arc=1.5pt,left=4pt,right=4pt,top=2pt,bottom=2pt,title={\scriptsize\bfseries #1},fonttitle=\color{gray},breakable}
\newtcolorbox{harnessbox}[1][]{colback=harnessbg,colframe=harnessborder,arc=1.5pt,left=4pt,right=4pt,top=2pt,bottom=2pt,title={\scriptsize\bfseries #1},fonttitle=\color{harnessborder!80!black},breakable}

\lstdefinestyle{rust}{language=C,backgroundcolor=\color{rustbg},frame=single,rulecolor=\color{rustborder},
  basicstyle=\ttfamily\fontsize{7.5}{9.5}\selectfont,keywordstyle=\color{accent}\bfseries,
  commentstyle=\color{gray}\itshape,stringstyle=\color{success},
  morekeywords={fn,pub,struct,enum,impl,trait,use,mod,crate,self,Self,let,mut,const,type,where,async,await,match,if,else,for,while,loop,return,break,Box,Vec,HashMap,BTreeMap,Option,Result,Ok,Err,Some,None,i64,u64,f64,u8,u16,u32,bool,String,usize,derive,Clone,Debug,Serialize,Deserialize},
  breaklines=true,showstringspaces=false,xleftmargin=2pt,xrightmargin=2pt,aboveskip=4pt,belowskip=4pt,tabsize=4}
\lstdefinestyle{yaml}{backgroundcolor=\color{codebg},frame=single,rulecolor=\color{codeborder},
  basicstyle=\ttfamily\fontsize{7.5}{9.5}\selectfont,breaklines=true,showstringspaces=false,
  xleftmargin=2pt,xrightmargin=2pt,aboveskip=4pt,belowskip=4pt}
\lstdefinestyle{sh}{backgroundcolor=\color{codebg},frame=single,rulecolor=\color{codeborder},
  basicstyle=\ttfamily\fontsize{7.5}{9.5}\selectfont,breaklines=true,showstringspaces=false,
  xleftmargin=2pt,xrightmargin=2pt,aboveskip=4pt,belowskip=4pt,
  morekeywords={cargo,isls,run,test,bench,validate,report,replay,ingest}}

\setlength{\parskip}{0.3em}\setlength{\parindent}{0pt}

\begin{document}

% ═══════════════════════════════════════════════════════════════════════════════
\begin{titlepage}\centering\vspace*{1cm}
{\Huge\bfseries\color{sectioncolor}ISLS Validation Harness}\\[0.4cm]
{\LARGE\color{accent}v1.0.0}\\[0.3cm]
{\large\color{accentdark}Deterministic Benchmark, Empirical Validation\\
\& Operational Metering Framework}

\vspace{0.8cm}
\begin{tcolorbox}[colback=thmbg,colframe=thmborder,width=0.92\textwidth,arc=3pt]
\small\centering
This document specifies the complete testing, benchmarking, and operational validation
layer that wraps the ISLS core. It is designed so that a single operator can:\\[2pt]
\textbf{(1)} launch the system against live or historical data with one command,\\
\textbf{(2)} collect all metrics automatically,\\
\textbf{(3)} read a single dashboard to determine system health,\\
\textbf{(4)} identify exactly which subsystem needs iteration.\\[2pt]
The harness itself is pure Rust, ships as two additional crates in the ISLS workspace,
and requires zero ongoing human attention during normal operation.
\end{tcolorbox}

\vspace{0.6cm}
\begin{tabular}{rl}
\textbf{Author:} & Sebastian Klemm \\
\textbf{Date:} & 12 March 2026 \\
\textbf{Depends on:} & ISLS Rust Implementation Specification v1.0.0 \\
\textbf{New crates:} & \texttt{isls-harness} (C10), \texttt{isls-cli} (C11) \\
\end{tabular}
\vfill
{\scriptsize\textcolor{gray}{Operator workload target: $<$5 minutes/week for a healthy system. All anomalies surface automatically.}}
\end{titlepage}

\tableofcontents\newpage

% ═══════════════════════════════════════════════════════════════════════════════
\section{Architecture}
% ═══════════════════════════════════════════════════════════════════════════════

Two new crates extend the ISLS workspace:

\begin{lstlisting}[style=yaml]
isls/
  crates/
    isls-harness/    # C10: benchmark runners, validators, metric collectors, reporters
    isls-cli/        # C11: single-binary operator interface (commands, dashboard, config)
\end{lstlisting}

\begin{specbox}[Design Rule]
The harness never touches ISLS internals via back-doors. It uses exclusively the public API of \texttt{isls-engine} and \texttt{isls-archive}. If a metric cannot be computed from the public API, the public API must be extended --- the harness must not circumvent encapsulation.
\end{specbox}

\subsection{Operator Interaction Model}

The entire system is operated through one binary: \texttt{isls}.

\begin{lstlisting}[style=sh]
isls init                          # generate default config + data dirs
isls ingest <source>               # attach a live or file-based data source
isls run                           # start the macro-step loop (foreground or daemon)
isls run --replay <run_descriptor> # deterministic replay from saved descriptor
isls bench                         # run full benchmark suite, emit report
isls validate                      # run validation suite against collected data
isls report                        # print current health dashboard to stdout
isls report --json                 # machine-readable metrics export
isls report --html > report.html   # self-contained HTML dashboard
isls status                        # one-line system health summary
\end{lstlisting}

\textbf{Target:} A healthy system requires the operator to run \texttt{isls status} and glance at the output. Everything else is automated.

% ═══════════════════════════════════════════════════════════════════════════════
\section{Metric Taxonomy}
\label{sec:metrics}
% ═══════════════════════════════════════════════════════════════════════════════

Every metric has: a unique ID, a human name, a computation method, a unit, a healthy range, and an alert threshold.

\subsection{Layer Health Metrics (M1--M5)}

One metric per ISLS layer, answering: ``Is this layer functioning correctly?''

\begin{longtable}{l l p{4.8cm} l l}
\toprule \textbf{ID} & \textbf{Layer} & \textbf{Computation} & \textbf{Healthy} & \textbf{Alert} \\
\midrule\endhead
M1 & L0 Observe & Ingestion rate: observations/sec over last 60s & $>0$ & $=0$ for $>30$s \\
M2 & L1 Persist & Graph growth: $\Delta|V| + \Delta|E|$ per macro-window & $\geq 0$ & Negative (data loss) \\
M3 & L2 Extract & Active constraints: $|S_{\text{active}}|$ after scan & $\geq 1$ & $0$ for $>1$h \\
M4 & L3 Consensus & Crystal rate: crystals committed per 24h & Profile-dep. & $0$ for $>$ configured silent period \\
M5 & L4 Morph & Mutation rate: mutations per 24h & $\geq 0$ & Negative invariant violation \\
\bottomrule
\end{longtable}

\subsection{Core Quality Metrics (M6--M14)}

\begin{longtable}{l l p{4.8cm} l l}
\toprule \textbf{ID} & \textbf{Name} & \textbf{Computation} & \textbf{Healthy} & \textbf{Alert} \\
\midrule\endhead
M6 & Replay fidelity & Run macro-step twice with saved $RD$; compare crystal digest sets & 100\% match & Any mismatch \\
M7 & Convergence rate & Wasserstein distance $\delta(w)$ between consecutive persistence diagrams & Decreasing & Increasing for $>5$ windows \\
M8 & Lattice stability & Mean free energy $\bar{F}$ of emitted crystals & $<0$ & $\geq 0$ \\
M9 & Gate selectivity & Fraction of macro-steps that pass Kairos gate & $0.01$--$0.30$ & $>0.5$ (too permissive) or $0$ (too strict) \\
M10 & Dual consensus MCI & Mean MCI of committed crystals & $\geq 0.90$ & $<0.80$ \\
M11 & PoR latency & Mean time from search$\to$commit in PoR FSM & Profile-dep. & $>2\times$ configured budget \\
M12 & Evidence integrity & Fraction of crystals passing \texttt{verify\_crystal} & 100\% & $<100\%$ \\
M13 & Operator version drift & Count of unmatched operator versions in archive & 0 & $>0$ \\
M14 & Storage efficiency & Cold-tier bytes per asset per month & $<3$ MB & $>10$ MB \\
\bottomrule
\end{longtable}

\subsection{Performance Metrics (M15--M19)}

\begin{longtable}{l l p{4.8cm} l l}
\toprule \textbf{ID} & \textbf{Name} & \textbf{Computation} & \textbf{Healthy} & \textbf{Alert} \\
\midrule\endhead
M15 & Macro-step latency & Wall-clock time per \texttt{macro\_step()} call & $<60$s & $>120$s \\
M16 & Memory footprint & RSS of process & $<4$ GB (2k assets) & $>8$ GB \\
M17 & Extraction throughput & Constraint candidates evaluated per second & $>100$ & $<10$ \\
M18 & Archive growth rate & Bytes appended to cold tier per 24h & Bounded & Unbounded growth \\
M19 & Carrier migration latency & Time from migration trigger to stable new phase & $<5$s & $>30$s \\
\bottomrule
\end{longtable}

\subsection{Empirical Domain Metrics (M20--M24)}

These measure how the system performs \emph{in the real world} when connected to a domain profile (e.g., crypto markets via DSHAE).

\begin{longtable}{l l p{4.8cm} l l}
\toprule \textbf{ID} & \textbf{Name} & \textbf{Computation} & \textbf{Healthy} & \textbf{Alert} \\
\midrule\endhead
M20 & Constraint hit rate & Fraction of emitted constraints that remain active 24h later & $>0.6$ & $<0.3$ \\
M21 & Crystal predictive value & For crystals signaling structure $X$: did $X$ persist? & $>0.5$ & $<0.3$ \\
M22 & Signal lead time & Time between crystal emission and observable market event & $>0$ (leading) & $\leq 0$ (lagging) \\
M23 & Basket quality lift & DSHAE Sharpe ratio with ISLS basket vs.\ random basket & Positive & Negative \\
M24 & Coverage growth & Number of tracked entities $|V|$ over time & Monotone increasing & Decrease (data loss) \\
\bottomrule
\end{longtable}

% ═══════════════════════════════════════════════════════════════════════════════
\section{Benchmark Suite}
\label{sec:bench}
% ═══════════════════════════════════════════════════════════════════════════════

Invoked via \texttt{isls bench}. All benchmarks are deterministic (fixed seeds, synthetic data). Results are appended to a local benchmark history for regression tracking.

\subsection{Benchmark Catalog}

\begin{longtable}{l l p{5cm} l}
\toprule \textbf{ID} & \textbf{Name} & \textbf{Procedure} & \textbf{Metric} \\
\midrule\endhead
B01 & Ingestion throughput & Ingest $10^6$ synthetic observations & obs/sec \\
B02 & Graph update scaling & Update graph with $N=100,500,2000,5000$ vertices & $\mu$s per update vs.\ $N$ \\
B03 & Extraction scaling & Run \texttt{inverse\_weave} on $N=50,200,1000$ point clouds & ms per scan vs.\ $N$ \\
B04 & Cascade contraction & Apply DK$\to$WF$\to$TS$\to$PC on synthetic 5D cloud & Contraction ratio, ms \\
B05 & Dual consensus overhead & Run primal+dual paths; measure wall time vs.\ single & Overhead \% \\
B06 & Crystal serialization & Serialize+hash $10^4$ crystals & $\mu$s per crystal \\
B07 & Replay speed & Replay 1000 macro-steps from saved $RD$ & steps/sec \\
B08 & Evidence verification & Verify $10^4$ crystals with $\bar{n}=50$ evidence entries each & $\mu$s per verify \\
B09 & Memory scaling & Run engine with $N=100\to5000$ entities; measure RSS & MB vs.\ $N$ \\
B10 & Full macro-step & End-to-end \texttt{macro\_step()} on reference dataset & ms, peak RSS \\
\bottomrule
\end{longtable}

\subsection{Regression Tracking}

\begin{lstlisting}[style=rust]
/// Benchmark result record
pub struct BenchResult {
    pub bench_id: String,
    pub timestamp: chrono::DateTime<chrono::Utc>,
    pub git_commit: String,
    pub metric_name: String,
    pub metric_value: f64,
    pub metric_unit: String,
    pub params: BTreeMap<String, String>,
}

/// Regression check: current vs. last 5 runs
pub fn check_regression(current: &BenchResult, history: &[BenchResult]) -> RegressionVerdict {
    let mean = history.iter().map(|h| h.metric_value).sum::<f64>() / history.len() as f64;
    let std = /* ... */;
    if current.metric_value > mean + 2.0 * std { RegressionVerdict::Regression }
    else if current.metric_value < mean - 2.0 * std { RegressionVerdict::Improvement }
    else { RegressionVerdict::Stable }
}
\end{lstlisting}

Output: \texttt{bench\_history.jsonl} — one JSON line per benchmark run, enabling time-series analysis.

% ═══════════════════════════════════════════════════════════════════════════════
\section{Validation Suite}
\label{sec:validate}
% ═══════════════════════════════════════════════════════════════════════════════

Invoked via \texttt{isls validate}. Runs against \emph{collected real data} (not synthetic). Answers: ``Is the system producing correct and useful results in practice?''

\subsection{Validation Levels}

\begin{specbox}[Three Validation Levels]
\begin{description}[nosep,leftmargin=1.5cm,style=sameline]
\item[V-Formal] Invariant checks: all 20 ISLS invariants verified on every committed crystal in the archive. Deterministic, runs offline. \texttt{isls validate --formal}
\item[V-Retro] Retrospective accuracy: compare past crystal predictions against what actually happened. Requires $\geq 7$ days of operational data. \texttt{isls validate --retro}
\item[V-Live] Live monitoring: continuous metric collection during \texttt{isls run}. Alerts on anomaly. Always-on, no separate command.
\end{description}
\end{specbox}

\subsection{Formal Validation (V-Formal)}

Iterate over every crystal in the archive and verify:

\begin{lstlisting}[style=rust]
pub fn validate_formal(archive: &Archive, pinned: &BTreeMap<String,String>)
    -> FormalReport
{
    let mut report = FormalReport::new();
    for crystal in archive.all_crystals() {
        // V-F1: Content address integrity
        report.check("content_address", crystal.crystal_id == content_address(crystal));
        // V-F2: Evidence chain integrity
        report.check("evidence_chain", verify_evidence_chain(&crystal.evidence_chain).is_ok());
        // V-F3: Operator version match
        report.check("operator_versions", verify_operator_versions(crystal, pinned).is_ok());
        // V-F4: Gate values above thresholds
        report.check("gate_kairos", crystal.commit_proof.gate_values.kairos);
        // V-F5: Dual consensus
        let cr = &crystal.commit_proof.consensus_result;
        report.check("dual_consensus",
            cr.primal_score >= cr.threshold && cr.dual_score >= cr.threshold
            && cr.mci >= 0.80);
        // V-F6: PoR trace monotonicity
        report.check("por_trace", verify_por_trace(&crystal.commit_proof.por_trace));
        // V-F7: Free energy < 0
        report.check("free_energy", crystal.free_energy < 0.0);
        // V-F8: Append-only (no past crystal altered)
        report.check("immutability", !archive.was_modified(&crystal.crystal_id));
    }
    report
}
\end{lstlisting}

\subsection{Retrospective Validation (V-Retro)}

For each crystal that made a structural claim (e.g., ``constraint $\mathfrak{g}_k$ binds assets $A$ and $B$''), check whether the claim held true in subsequent data:

\begin{lstlisting}[style=rust]
pub fn validate_retro(archive: &Archive, graph: &PersistentGraph,
    horizon: Duration) -> RetroReport
{
    let mut report = RetroReport::new();
    for crystal in archive.crystals_older_than(horizon) {
        for constraint in &crystal.constraint_program {
            // Re-evaluate constraint on data AFTER crystal emission
            let post_window = TimeWindow::after(crystal.created_at, horizon);
            let coverage = evaluate_constraint_coverage(constraint, graph, &post_window);
            report.add(RetroEntry {
                crystal_id: crystal.crystal_id,
                constraint_id: constraint.id,
                predicted_coverage: constraint.coverage,
                actual_coverage: coverage,
                still_active: coverage >= 0.5, // did the constraint survive?
            });
        }
    }
    // Aggregate: hit rate, mean coverage drift, false positive rate
    report.compute_aggregates();
    report
}
\end{lstlisting}

\textbf{Key output:} Constraint hit rate (M20), Crystal predictive value (M21), Signal lead time (M22).

\subsection{Live Validation (V-Live)}

Runs as a background task inside \texttt{isls run}. On every macro-step:

\begin{enumerate}[nosep]
\item Compute all 24 metrics (M1--M24).
\item Append to \texttt{metrics.jsonl} (one JSON line per tick).
\item Check each metric against alert threshold.
\item If any alert fires: log \texttt{WARN} to structured log + write to \texttt{alerts.jsonl}.
\item Every $N$ macro-steps (configurable, default 100): compute rolling aggregates and append to \texttt{summary.jsonl}.
\end{enumerate}

% ═══════════════════════════════════════════════════════════════════════════════
\section{Dashboard and Reporting}
\label{sec:dashboard}
% ═══════════════════════════════════════════════════════════════════════════════

\subsection{One-Line Status}

\begin{cmdbox}[\texttt{isls status}]
\begin{lstlisting}[style=sh]
$ isls status
ISLS v1.0.0 | UP 14d 3h | Entities: 847 | Edges: 12,403 | Crystals: 23
L0: OK (42 obs/s) | L1: OK (+12 edges/h) | L2: OK (31 active constraints)
L3: OK (1.3 crystals/day, MCI=0.94) | L4: OK (0 mutations pending)
Replay: PASS | Free Energy: -0.38 | Gate sel.: 0.07 | Storage: 2.1 GB
Health: GREEN
\end{lstlisting}
\end{cmdbox}

Color codes: \textcolor{success}{\textbf{GREEN}} = all metrics healthy. \textcolor{warn}{\textbf{YELLOW}} = $\geq 1$ metric in warning range. \textcolor{danger}{\textbf{RED}} = $\geq 1$ metric in alert.

\subsection{Full Report}

\texttt{isls report} outputs a structured report with sections:

\begin{enumerate}[nosep]
\item \textbf{System Overview:} uptime, entity count, crystal count, storage usage.
\item \textbf{Layer Health:} M1--M5 with traffic-light indicators.
\item \textbf{Core Quality:} M6--M14 with sparkline trend (last 30 data points).
\item \textbf{Performance:} M15--M19 with p50/p95/p99 latencies.
\item \textbf{Empirical:} M20--M24 with rolling 7-day and 30-day aggregates.
\item \textbf{Alerts:} Active alerts with timestamp, metric ID, and current value.
\item \textbf{Iteration Guidance:} Ranked list of subsystems needing attention (see \S\ref{sec:iteration}).
\end{enumerate}

\texttt{isls report --html} emits a self-contained HTML file with embedded CSS and SVG charts (no JavaScript dependencies, no external fetches). Viewable in any browser.

\texttt{isls report --json} emits machine-readable JSON for programmatic consumption.

% ═══════════════════════════════════════════════════════════════════════════════
\section{Iteration Guidance Engine}
\label{sec:iteration}
% ═══════════════════════════════════════════════════════════════════════════════

The harness does not just report metrics --- it tells the operator \emph{what to fix next}.

\subsection{Diagnosis Rules}

\begin{longtable}{l p{3.5cm} p{5.5cm} l}
\toprule \textbf{Trigger} & \textbf{Symptom} & \textbf{Diagnosis + Action} & \textbf{Priority} \\
\midrule\endhead
M6 $<100\%$ & Replay mismatch & Non-determinism in operator or serialization. Run \texttt{isls replay --diff} to identify first diverging macro-step. Check for \texttt{HashMap} usage, floating-point platform difference, or operator version drift. & \textcolor{danger}{P0} \\
M8 $\geq 0$ & Positive free energy & Constraints too weak or temperature too high. Lower \texttt{extraction.alpha\_min} or recalibrate temperature $c_T$. & \textcolor{warn}{P1} \\
M9 $>0.5$ & Gate too permissive & Thresholds too low. Increase \texttt{thresholds.k} (crystal stability). Review whether noise is being crystallized. & \textcolor{warn}{P1} \\
M9 $=0$ for $>48$h & Gate too strict & Thresholds too high for current data regime. Lower \texttt{thresholds.d} and \texttt{thresholds.q} by 10\% increments. & \textcolor{warn}{P1} \\
M10 $<0.80$ & Low MCI & Primal and dual paths diverge. Check operator cascade ordering. May indicate that DK contraction rate is too aggressive --- reduce \texttt{dk\_alpha}. & \textcolor{warn}{P1} \\
M15 $>120$s & Slow macro-step & Profile to find bottleneck: likely L2 extraction. Reduce \texttt{extraction.max\_iterations} or \texttt{press\_top\_k}. Consider sparser constraint library. & \textcolor{warn}{P2} \\
M20 $<0.3$ & Low constraint hit rate & Constraints are transient noise, not structural. Increase \texttt{extraction.bond\_strength\_min} so only durable constraints enter programs. & \textcolor{warn}{P1} \\
M21 $<0.3$ & Low predictive value & Crystals do not predict future structure. Likely cause: macro-window too short. Increase \texttt{extraction.window\_hours}. & \textcolor{warn}{P1} \\
M22 $\leq 0$ & Lagging signals & ISLS detects structure \emph{after} the market has moved. Reduce \texttt{por\_T\_min} and \texttt{por\_T\_stable} to allow faster PoR traversal. & \textcolor{warn}{P2} \\
M12 $<100\%$ & Evidence corruption & Storage integrity failure. Run \texttt{isls validate --formal}. Identify corrupted tier segment. Re-derive from cold evidence or isolate. & \textcolor{danger}{P0} \\
\bottomrule
\end{longtable}

\subsection{Automated Iteration Report}

\begin{lstlisting}[style=rust]
pub struct IterationItem {
    pub priority: Priority,       // P0, P1, P2
    pub metric_id: String,        // e.g. "M8"
    pub subsystem: String,        // e.g. "isls-extract"
    pub symptom: String,          // human-readable
    pub diagnosis: String,        // what's likely wrong
    pub action: String,           // what to change
    pub config_key: Option<String>, // which config field to adjust
    pub suggested_value: Option<String>, // suggested new value
}

pub fn generate_iteration_guidance(metrics: &MetricSnapshot, config: &Config)
    -> Vec<IterationItem>
{
    let mut items = Vec::new();
    if metrics.m6_replay_fidelity < 1.0 {
        items.push(IterationItem {
            priority: Priority::P0,
            metric_id: "M6".into(),
            subsystem: "isls-engine".into(),
            symptom: "Replay mismatch detected".into(),
            diagnosis: "Non-determinism in operator or serialization path".into(),
            action: "Run `isls replay --diff` to find first diverging step".into(),
            config_key: None,
            suggested_value: None,
        });
    }
    // ... (one rule per row in the diagnosis table above)
    items.sort_by_key(|i| i.priority.clone());
    items
}
\end{lstlisting}

The output of \texttt{isls report} includes an \textbf{``Action Items''} section listing these in priority order. The operator reads top to bottom and addresses P0 first.

% ═══════════════════════════════════════════════════════════════════════════════
\section{Data Sources and Ingestion Adapters}
\label{sec:sources}
% ═══════════════════════════════════════════════════════════════════════════════

\subsection{Built-In Adapters}

\begin{tabular}{l l l}
\toprule \textbf{Adapter} & \textbf{Source} & \textbf{Invocation} \\
\midrule
\texttt{file-csv} & CSV files (historical OHLCV) & \texttt{isls ingest --adapter file-csv --path data/} \\
\texttt{file-jsonl} & JSONL observation logs & \texttt{isls ingest --adapter file-jsonl --path obs.jsonl} \\
\texttt{replay} & Saved RunDescriptor & \texttt{isls run --replay run\_desc.json} \\
\texttt{synthetic} & Deterministic synthetic generator & \texttt{isls ingest --adapter synthetic --entities 500} \\
\bottomrule
\end{tabular}

\texttt{synthetic} generates deterministic data with configurable correlation structure, regime switches, and known constraint programs --- enabling ground-truth validation (the harness \emph{knows} which constraints exist and can measure detection accuracy exactly).

\subsection{Extension Point for Live Feeds}

\begin{lstlisting}[style=rust]
/// Trait for live data source adapters (domain-specific, out of ISLS core scope)
pub trait LiveAdapter: ObservationAdapter {
    fn connect(&mut self) -> Result<()>;
    fn poll(&mut self) -> Result<Vec<Vec<u8>>>;  // raw observation bytes
    fn disconnect(&mut self);
}
// A crypto-market adapter implementing this trait would connect to exchange WebSockets.
// The harness provides the framework; the domain-specific adapter is a plugin.
\end{lstlisting}

% ═══════════════════════════════════════════════════════════════════════════════
\section{Operational Modes}
\label{sec:modes}
% ═══════════════════════════════════════════════════════════════════════════════

\begin{specbox}[Four Operational Modes]
\begin{description}[nosep,leftmargin=2.5cm,style=sameline]
\item[Bench] Synthetic data, deterministic seeds. Measures performance. No side effects. \\ \texttt{isls bench}
\item[Validate] Real archived data, offline. Measures correctness and predictive quality. \\ \texttt{isls validate}
\item[Shadow] Live data ingestion, full pipeline, but \textbf{no downstream action}. Crystals are emitted and archived but not forwarded to DSHAE. Enables safe burn-in. \\ \texttt{isls run --mode shadow}
\item[Live] Full operation. Crystals forwarded to downstream consumers. Metrics collected. \\ \texttt{isls run --mode live}
\end{description}
\end{specbox}

\textbf{Recommended deployment sequence:}
\begin{enumerate}[nosep]
\item Run \texttt{isls bench} after every code change. Ensure no regressions.
\item Deploy in \texttt{shadow} mode for $\geq 14$ days. Collect metrics.
\item Run \texttt{isls validate --retro} after 14 days. Check M20--M24.
\item If all GREEN: switch to \texttt{live} mode.
\item Run \texttt{isls report} weekly. Address iteration items top-down.
\end{enumerate}

% ═══════════════════════════════════════════════════════════════════════════════
\section{File Layout and Persistence}
\label{sec:files}
% ═══════════════════════════════════════════════════════════════════════════════

\begin{lstlisting}[style=yaml]
~/.isls/
  config.json              # master config (ISLS + harness)
  data/
    hot/                   # hot-tier observation cache (7 days)
    warm/                  # warm-tier compressed archive (90 days)
    cold/                  # cold-tier indefinite evidence
    crystals/              # committed crystal archive (one file per crystal)
  metrics/
    metrics.jsonl          # per-tick metric log (M1-M24)
    alerts.jsonl           # fired alerts
    summary.jsonl          # rolling aggregates (every 100 ticks)
    bench_history.jsonl    # benchmark regression history
  reports/
    latest.html            # most recent HTML report
    latest.json            # most recent JSON report
  replay/
    *.run_descriptor.json  # saved run descriptors for replay
\end{lstlisting}

% ═══════════════════════════════════════════════════════════════════════════════
\section{Configuration Extension}
\label{sec:config}
% ═══════════════════════════════════════════════════════════════════════════════

The harness adds a \texttt{harness} section to the ISLS config:

\begin{lstlisting}[style=yaml]
harness:
  metrics_interval_steps: 1        # collect metrics every N macro-steps
  summary_interval_steps: 100      # rolling summary every N steps
  alert_cooldown_seconds: 300      # min time between repeated alerts
  bench:
    synthetic_entities: 500
    synthetic_windows: 1000
    seed: 42
  validate:
    retro_horizon_days: 7
    formal_on_startup: true        # run formal validation on every startup
  report:
    auto_generate_hours: 24        # auto-generate HTML report every N hours
    retention_days: 90             # keep old reports for N days
  iteration:
    enabled: true                  # auto-generate iteration guidance
    p0_halt: false                 # if true, halt engine on P0 alert (safety mode)
\end{lstlisting}

% ═══════════════════════════════════════════════════════════════════════════════
\section{Synthetic Ground-Truth Generator}
\label{sec:synthetic}
% ═══════════════════════════════════════════════════════════════════════════════

The most powerful validation tool: a deterministic data generator that embeds \emph{known} constraint programs into synthetic data, then measures ISLS's ability to recover them.

\begin{lstlisting}[style=rust]
pub struct SyntheticScenario {
    pub name: String,
    pub entities: usize,
    pub windows: usize,
    pub seed: u64,
    pub planted_constraints: Vec<PlantedConstraint>,
    pub regime_switches: Vec<(usize, RegimeType)>,  // (window, new_regime)
}

pub struct PlantedConstraint {
    pub template: ConstraintTemplate,
    pub parameters: BTreeMap<String, f64>,
    pub active_from: usize,   // window number
    pub active_until: usize,
    pub strength: f64,        // how strongly the constraint is enforced in data
}

/// Generate synthetic observations with planted constraints
pub fn generate(scenario: &SyntheticScenario) -> Vec<Vec<Observation>> { /* ... */ }

/// Score ISLS's recovery of planted constraints
pub fn score_recovery(planted: &[PlantedConstraint], discovered: &ConstraintProgram)
    -> RecoveryScore
{
    // Precision: fraction of discovered constraints that match a planted one
    // Recall: fraction of planted constraints that were discovered
    // F1: harmonic mean
    // Parameter accuracy: mean |param_discovered - param_planted| for matched constraints
    RecoveryScore { precision, recall, f1, param_accuracy }
}
\end{lstlisting}

\textbf{Reference scenarios:}
\begin{enumerate}[nosep]
\item \textbf{S-Basic:} 50 entities, 3 planted Band constraints, no regime switches.
\item \textbf{S-Regime:} 200 entities, 5 constraints, 2 regime switches (calm$\to$volatile$\to$calm).
\item \textbf{S-Causal:} 100 entities, 3 Granger causal chains $A\to B\to C$, planted with 2h lag.
\item \textbf{S-Break:} 200 entities, 4 constraints active for 500 windows then breaking.
\item \textbf{S-Scale:} 2000 entities, 20 constraints, full regime complexity.
\end{enumerate}

\texttt{isls bench} runs all five scenarios and reports precision/recall/F1 for each.

% ═══════════════════════════════════════════════════════════════════════════════
\section{Implementation Contract for Claude Code}
\label{sec:contract}
% ═══════════════════════════════════════════════════════════════════════════════

\begin{harnessbox}[Build Contract]
\small
\begin{enumerate}[nosep]
\item Create crate \texttt{isls-harness} with modules: \texttt{metrics}, \texttt{bench}, \texttt{validate}, \texttt{synthetic}, \texttt{report}, \texttt{iterate}.
\item Create crate \texttt{isls-cli} with binary \texttt{isls} implementing all commands from \S1.
\item Implement all 24 metrics (M1--M24) with collection, alert, and persistence.
\item Implement all 10 benchmarks (B01--B10) with regression tracking.
\item Implement V-Formal, V-Retro, V-Live validation levels.
\item Implement the 10 diagnosis rules from \S\ref{sec:iteration} and the iteration guidance generator.
\item Implement the synthetic ground-truth generator with all 5 reference scenarios.
\item Implement \texttt{isls report --html} with self-contained HTML (inline CSS, SVG sparklines, no external deps).
\item All tests: \texttt{cargo test --workspace} must pass.
\item The binary \texttt{isls status} must produce output matching the format in \S5.1 within 1 second.
\end{enumerate}
\end{harnessbox}

\vspace{0.5cm}
\begin{center}
\rule{0.6\textwidth}{0.3pt}\\[0.3cm]
{\small End of specification — ISLS Validation Harness v1.0.0}\\[2pt]
{\scriptsize 24 metrics. 10 benchmarks. 3 validation levels. 10 diagnosis rules. 5 synthetic scenarios.\\One binary. One glance.}
\end{center}

\end{document}
